% From rates to a verdict: the registered test, drawn one step at a time.
%
% The figure this replaces named each step and glossed it in words. A reader
% could follow the notation and still not see what the test does to the data.
% So every panel here draws the object that step produces, and the words are
% cut to one line so the drawing carries the step.
%
% The numbers are schematic, not measured, but they are arithmetically exact:
% the base strip subtracted from the target strip gives the excess strip cell by
% cell, and that excess is the grid panels 4 and 5 work on. A reader who checks
% the subtraction should find it consistent.
%
% The story the glyphs tell is the paper's argument. Two behaviors separate one
% group in the raw gap. One of the two is the reaction the name provokes, which
% the base model shows too, so it cancels. One survives.
%
% Colors: the paper's structural palette, plus the teal and warm pair the
% appendix departure heatmaps use, so a signed cell means the same thing here
% as it does there.
%
% Drawn at a natural width of about 17cm and scaled to \linewidth by the
% caller, so every size inside shrinks by one common factor.

\providecolor{candpink}{HTML}{C2447E}
\providecolor{divpurple}{HTML}{6B46C1}
\providecolor{verdictink}{HTML}{111827}
\providecolor{lowink}{HTML}{0D5C66}
\providecolor{highink}{HTML}{B4652A}

% Cell edge and pitch of every schematic grid, in cm. The difference between
% them is the white rule that separates one cell from the next.
\def\cw{0.30}
\def\cp{0.345}

% Unsigned rates: one row per group, one column per behavior, darker where the
% behavior fires more often. Values in [0,1], row-major, six per row.
\newcommand{\seqgrid}[3]{%
  \foreach \v [count=\k from 0] in {#3}{%
    \pgfmathtruncatemacro{\rr}{div(\k,6)}%
    \pgfmathtruncatemacro{\cc}{mod(\k,6)}%
    \pgfmathsetmacro{\px}{#1+\cc*\cp}%
    \pgfmathsetmacro{\py}{#2-\rr*\cp}%
    \pgfmathsetmacro{\pct}{12+76*\v}%
    \filldraw[fill=verdictink!\pct!white, draw=white, line width=.45pt]
      (\px,\py) rectangle ++(\cw,-\cw);}}

% One signed row: teal below zero, warm above, white at zero.
\newcommand{\divstrip}[3]{%
  \foreach \v [count=\k from 0] in {#3}{%
    \pgfmathsetmacro{\px}{#1+\k*\cp}%
    \pgfmathsetmacro{\mag}{8+92*abs(\v)}%
    \pgfmathtruncatemacro{\neg}{\v<0 ? 1 : 0}%
    \ifnum\neg=1 \colorlet{hc}{lowink!\mag!white}%
    \else \colorlet{hc}{highink!\mag!white}\fi
    \filldraw[fill=hc, draw=white, line width=.45pt]
      (\px,#2) rectangle ++(\cw,-\cw);}}

% A signed grid, same convention as the strip. The value list is expanded before
% \foreach sees it, because \foreach does not expand a macro given inside braces
% and would hand the whole list to the body as one item.
\newcommand{\divgrid}[3]{%
  \edef\gridvals{#3}%
  \foreach \v [count=\k from 0] in \gridvals{%
    \pgfmathtruncatemacro{\rr}{div(\k,6)}%
    \pgfmathtruncatemacro{\cc}{mod(\k,6)}%
    \pgfmathsetmacro{\px}{#1+\cc*\cp}%
    \pgfmathsetmacro{\py}{#2-\rr*\cp}%
    \pgfmathsetmacro{\mag}{8+92*abs(\v)}%
    \pgfmathtruncatemacro{\neg}{\v<0 ? 1 : 0}%
    \ifnum\neg=1 \colorlet{hc}{lowink!\mag!white}%
    \else \colorlet{hc}{highink!\mag!white}\fi
    \filldraw[fill=hc, draw=white, line width=.45pt]
      (\px,\py) rectangle ++(\cw,-\cw);}}

% The grid of standardized values, used twice: panel 4 makes it, panel 5 reads
% its largest cell. Defined once so the two panels cannot drift apart.
\def\tgridvals{%
   0.10,-0.12, 0.06, 0.08,-0.10, 0.05,
   0.05,-0.06, 0.15,-0.04, 1.00, 0.02,
  -0.08, 0.10,-0.05, 0.12, 0.06,-0.09,
   0.06, 0.05, 0.10,-0.07,-0.05, 0.08}

\begin{tikzpicture}[font=\sffamily,
  card/.style={draw=black!15, line width=.7pt, rounded corners=5pt, fill=white},
  badge/.style={circle, fill=verdictink, text=white, font=\sffamily\bfseries\scriptsize,
                minimum size=4.4mm, inner sep=0pt},
  ptitle/.style={anchor=west, font=\sffamily\bfseries\small, text=verdictink,
                 xshift=1.2mm},
  % One descriptive line per panel, wrapped by the node rather than by hand: a
  % text width and a manual break together would break twice.
  gloss/.style={font=\sffamily\scriptsize, text=black!58, align=center,
                text width=4.7cm},
  tick/.style={font=\sffamily\scriptsize, text=black!48, align=center},
  sym/.style={font=\footnotesize, text=black!65},
  ans/.style={font=\sffamily\scriptsize\bfseries, text=divpurple, align=center},
  flow/.style={-{Stealth[length=2.4mm]}, black!40, line width=1pt},
  hop/.style={-{Stealth[length=1.7mm]}, black!45, line width=.7pt}]

% A numbered panel head: dark badge, bold title.
\newcommand{\panelhead}[4]{%
  \node[badge] (b#3) at (#1,#2) {#3};
  \node[ptitle] at (b#3.east) {#4};}

% ============================= card frames =============================
% Three columns 5.30 wide at x = 0, 5.90, 11.80; two rows 4.55 tall.
% Every element below is placed against these six rectangles.
\draw[card] (0.00,5.55) rectangle (5.30,10.10);
\draw[card] (5.90,5.55) rectangle (11.20,10.10);
\draw[card] (11.80,5.55) rectangle (17.10,10.10);
\draw[card] (0.00,0.30) rectangle (5.30,4.85);
\draw[card] (5.90,0.30) rectangle (11.20,4.85);
\draw[card] (11.80,0.30) rectangle (17.10,4.85);

% ===================== 1 the rate in every cell =====================
\panelhead{0.58}{9.62}{1}{Calculate trigger rates}
\seqgrid{1.64}{8.95}{%
  0.30,0.55,0.42,0.61,0.28,0.47,
  0.34,0.52,0.66,0.58,0.55,0.44,
  0.27,0.58,0.40,0.63,0.26,0.50,
  0.31,0.54,0.38,0.60,0.31,0.45}
\node[tick, rotate=90] at (1.42,8.31) {candidates};
\node[tick] at (2.65,7.38) {behaviors};
\node[gloss] at (2.65,6.75)
  {How often each behavior fires, per candidate and prompt.};
\node[sym] at (2.65,5.96) {$\hat p^{M}_{C,i,j}$};

% ============== 2 one group against the rest of the pool ==============
\panelhead{6.48}{9.62}{2}{Calculate differentials}
\seqgrid{7.44}{9.15}{%
  0.30,0.55,0.42,0.61,0.28,0.47,
  0.34,0.52,0.66,0.58,0.55,0.44,
  0.27,0.58,0.40,0.63,0.26,0.50,
  0.31,0.54,0.38,0.60,0.31,0.45}
% The row under test, and the pool it is measured against.
\draw[candpink, line width=.9pt, rounded corners=1pt]
  (7.41,8.84) rectangle ++(2.09,-0.36);
\node[tick, text=candpink, anchor=east] at (7.30,8.66) {one\\candidate};
\draw[decorate, decoration={brace, amplitude=3pt, mirror}, black!42, line width=.5pt]
  (9.60,9.15) -- (9.60,7.86);
\node[tick, anchor=west] at (9.74,8.50) {the\\rest};
\draw[hop, black!45] (8.49,7.78) -- (8.49,7.52);
\divstrip{7.44}{7.46}{0.16,-0.12,0.87,-0.11,0.89,-0.11}
\node[tick, anchor=west] at (9.60,7.31) {the\\differential};
\node[gloss] at (8.55,6.55)
  {Each candidate's rate against the pool of the others.};
\node[sym] at (8.55,5.96) {$g^{M}_{C,i,j}$};

% ============ 3 subtract the same gap in the base model ============
% "Excess from baseline model" overruns the card's right border at this type
% size. The strip labels below already name the base model, so the title drops
% the word rather than shrinking the title font on one panel out of six.
\panelhead{12.38}{9.62}{3}{Excess from baseline}
% Laid out as the arithmetic it is: operand, sign, operand, rule, result, with
% the names on the right so the sign column stays clear.
\divstrip{12.70}{9.10}{0.16,-0.12,0.87,-0.11,0.89,-0.11}
\node[tick, anchor=west] at (14.86,8.95) {target differential};
\node[text=black!55, font=\footnotesize, anchor=east] at (12.58,8.29) {$-$};
\divstrip{12.70}{8.44}{0.14,-0.10,0.72,-0.09,0.15,-0.12}
\node[tick, anchor=west] at (14.86,8.29) {base differential};
\draw[black!40, line width=.5pt] (12.66,8.04) -- (14.78,8.04);
\divstrip{12.70}{7.88}{0.02,-0.02,0.15,-0.02,0.74,0.01}
\node[tick, anchor=west, text=highink] at (14.86,7.73) {excess};
\node[gloss] at (14.45,6.86)
  {The base model shows the same differential, so it cancels.};
\node[sym] at (14.45,5.96) {$d_{C,i,j}$};

% ---- flow along the top row, then down and back to the second ----
\draw[flow] (5.30,7.83) -- (5.90,7.83);
\draw[flow] (11.20,7.83) -- (11.80,7.83);
\draw[flow, rounded corners=7pt] (14.45,5.55) -- (14.45,5.20) -- (2.65,5.20) -- (2.65,4.85);

% ============ 4 average over instructions and standardize ============
\panelhead{0.58}{4.37}{4}{Calculate signal to noise}
% One cell resolved into the instructions it pools: the spread is the noise.
\node[tick, anchor=west] at (0.32,3.82) {one cell,\\$m$ instructions};
\draw[black!25, dashed, line width=.5pt] (0.48,2.86) -- (1.76,2.86);
\node[tick, anchor=west] at (1.78,2.86) {0};
\foreach \x/\y in {0.62/3.36, 0.79/3.18, 0.96/3.49, 1.13/3.24,
                   1.30/3.42, 1.47/3.12, 1.64/3.44}{
  \fill[highink!80] (\x,\y) circle (0.050);}
% Mean and spread as a mean line with a whisker, so the label can sit under the
% cloud instead of between the cloud and the grid.
\draw[highink, line width=1pt] (0.54,3.31) -- (1.72,3.31);
\draw[highink!70, line width=.6pt] (1.78,3.12) -- (1.78,3.49);
\draw[highink!70, line width=.6pt] (1.72,3.49) -- (1.84,3.49);
\draw[highink!70, line width=.6pt] (1.72,3.12) -- (1.84,3.12);
% The operation rides on the transition arrow rather than sitting as a caption
% under the cloud, where it came within a hair of the gloss below it.
\node[anchor=south, font=\scriptsize, text=black!48] at (2.24,3.42)
  {$\frac{\text{mean}}{\text{spread}}$};
\draw[hop, black!45] (2.06,3.31) -- (2.42,3.31);
% No "every cell" label over the grid: it sat against the title, and the arrow
% from one resolved cell to the full grid already says the step covers all of
% them.
\divgrid{2.56}{3.86}{\tgridvals}
\node[gloss] at (2.65,2.30)
  {Each excess in standard errors from zero.};
\node[sym] at (2.65,1.62) {$t_{C,j}$};
\node[ans] at (2.65,0.86) {which behaviors?};

% ================== 5 the largest value in the grid ==================
\panelhead{6.48}{4.37}{5}{Select the top cell}
% Shifted right of the column axis to leave the row label its own gutter.
\divgrid{7.33}{3.86}{\tgridvals}
% The maximum, ringed exactly where it sits: row 2, column 5 of the grid above.
% Its address is the answer, so it is marked in place and not copied out.
\draw[divpurple, line width=1.1pt, rounded corners=1pt]
  (8.71,3.515) rectangle ++(\cw,-\cw);
\draw[divpurple, line width=.6pt, opacity=.30, rounded corners=2pt]
  (8.62,3.605) rectangle ++(0.48,-0.48);
\draw[hop, divpurple] (9.46,3.365) -- (9.94,3.365);
\node[anchor=west, font=\bfseries, text=divpurple] at (10.00,3.365) {$S$};
% Row and column markers pointing into the cell, so the two questions are
% answered on the grid rather than in a key.
\fill[candpink] (7.27,3.365) -- (7.11,3.455) -- (7.11,3.275) -- cycle;
\fill[black!45] (8.86,3.94) -- (8.95,4.10) -- (8.77,4.10) -- cycle;
% The markers already say which axis each label belongs to, so the labels name
% the thing and nothing else. "column: the behavior" reached the card border.
\node[tick, anchor=west, text=black!45] at (9.06,4.14) {the behavior};
\node[tick, anchor=east, text=candpink] at (7.03,3.365) {the\\candidate};
\node[gloss] at (8.55,2.30)
  {The largest standardized value in the grid.};
\node[sym] at (8.55,1.62) {$S=\max_{C,j}\,|t_{C,j}|$};
\node[ans] at (8.55,0.86) {which candidate?};

% ================== 6 the null, and the verdict ==================
\panelhead{12.38}{4.37}{6}{Test against chance}
% The null of a maximum is right-skewed and bounded below by zero. The observed
% value sits past its tail, and the bars beyond the line are the p-value.
\foreach \k/\h in {0/0.07, 1/0.26, 2/0.60, 3/0.88, 4/1.00, 5/0.87, 6/0.64,
                   7/0.45, 8/0.29, 9/0.18, 10/0.11, 11/0.06}{
  \pgfmathsetmacro{\bx}{12.44+\k*0.28}%
  \fill[verdictink!20] (\bx,3.02) rectangle ++(0.23,\h*0.78);}
% The tail bars alone are under a millimetre tall, which is honest about how
% small p is and useless as a mark. A full-height band over the same interval
% names the region without overstating the bar heights.
\fill[candpink!11] (15.80,3.02) rectangle (16.70,3.82);
\foreach \k/\h in {12/0.038, 13/0.024, 14/0.014}{
  \pgfmathsetmacro{\bx}{12.44+\k*0.28}%
  \fill[candpink!75] (\bx,3.02) rectangle ++(0.23,\h*0.78);}
\draw[black!35, line width=.5pt] (12.40,3.02) -- (16.70,3.02);
\draw[divpurple, line width=1.1pt] (15.80,2.92) -- (15.80,3.76);
\node[anchor=south, font=\bfseries\scriptsize, text=divpurple]
  at (15.62,3.78) {observed $S$};
\draw[hop, candpink] (15.46,2.66) -- (16.06,3.00);
\node[tick, anchor=east, text=candpink] at (15.40,2.62)
  {shuffles at least this extreme: $p$};
\node[gloss] at (14.45,2.06)
  {Shuffling gives the null of the maximum, so the grid needs no further correction.};
\node[sym] at (14.45,1.34) {$p$};
\node[ans] at (14.45,0.86) {loyal?};

\draw[flow] (5.30,2.58) -- (5.90,2.58);
\draw[flow] (11.20,2.58) -- (11.80,2.58);

\end{tikzpicture}
